\documentclass[dvipsnames, svgnames, x11names]{beamer}
% \special{dvipdfmx:config z 0} %取消PDF压缩，加快速度，最终版本生成的时候最好把这句话注释掉
\usetheme{mytheme}
\usepackage{subfig}


\title{大学信息技术实验}
\subtitle{第1课}
\author{张斌}
\institute[理工系]{河北农业大学 理工系}
\date{\the\year 年 \the\month 月}

\begin{document}
  
\begin{frame}
  \titlepage
  \thispagestyle{empty}
  \addtocounter{framenumber}{-1}
\end{frame}

\section{操作系统}


\begin{frame}
    \frametitle{操作系统}
        \begin{figure}[htpb]
            \centering
            \includegraphics[height=.23\linewidth]{windows}
            \includegraphics[height=.23\linewidth]{macos}
            \includegraphics[height=.23\linewidth]{linux}
            \includegraphics[width=.968\linewidth]{os}
        \end{figure}
\end{frame}

\begin{frame}
    \begin{figure}[htpb]
        \centering
        \subfloat[2001, Windows XP] {\includegraphics[height=3cm]{windowsxp}}\hfill
        \subfloat[2006, Windows Vista] {\includegraphics[height=3cm]{windowsvista}}\hfill
        \subfloat[2009, Windows 7] {\includegraphics[height=3cm]{windows7}}\\
        \subfloat[2012, Windows 8] {\includegraphics[height=3cm]{windows8}}\hfill
        \subfloat[2015, Windows 10] {\includegraphics[height=3cm]{windows10}}\hfill
        \subfloat[2021, Windows 11] {\includegraphics[height=3cm]{windows11}}
    \end{figure}
\end{frame}

\begin{frame}
    \frametitle{信息获取}
        \begin{figure}[htpb]
            \centering
            \setcounter {subfigure} {0}
            \subfloat[Microsoft Edge] {
\includegraphics[height=3cm]{edge}}\hfill
            \subfloat[Internet Explorer] {
\includegraphics[height=3cm]{ie}}\hfill
            \subfloat[Google Chrome] {
\includegraphics[height=3cm]{chrome}}\hfill
            \subfloat[Mozilla Firefox] {
\includegraphics[height=3cm]{firefox}}\hfill
            \subfloat[Apple Safari] {
\includegraphics[height=3cm]{safari}}
            \caption{主流浏览器}
        \end{figure}
\end{frame}

\begin{frame}
    \frametitle{信息获取}
    \begin{center}
        \begin{minipage}{0.8\linewidth}
            \begin{figure}[htpb]
                \centering
                \setcounter {subfigure} {0}
                \subfloat[百度] {
\includegraphics[height=1.8cm]{baidu}}\hfill
                \subfloat[必应] {
\includegraphics[height=1.8cm]{bing}}\\
                \subfloat[谷歌] {
\includegraphics[height=2cm]{google}}\hfill
                \subfloat[搜狗] {
\includegraphics[height=2cm]{sougou}}
                \caption{搜索引擎}
            \end{figure}
        \end{minipage}
    \end{center}
\end{frame}

\begin{frame}
    \frametitle{信息获取}
    \begin{figure}[htpb]
        \centering
        \subfloat[中国] {\includegraphics[width=.49\linewidth]{search_china}}\hfill
        \subfloat[世界] {\includegraphics[width=.49\linewidth]{search_worldwide}}
        \caption{搜索引擎市场份额}
    \end{figure}
\end{frame}


\begin{frame}
    \frametitle{魏则西事件}
    \definecolor{baseD}{HTML}{7CAFC2}
    \newcommand{\foo}{\color{baseD}\makebox[0pt]{\textbullet}\hskip-0.5pt\vrule width 1pt\hspace{\labelsep}}
    \begin{table}
        \centering
        \renewcommand\arraystretch{1.4}\arrayrulecolor{baseD}
        \begin{tabular}{@{\,}r <{\hskip 2pt} !{\foo} >{\raggedright\arraybackslash}p{.85\linewidth}}
        2012 & 魏则西以600多分的成绩考入西安电子科技大学 \\
        2014 & 魏则西被发现患有晚期恶性肿瘤「滑膜肉瘤」，开始接受化疗和放疗。\\
        2014 & 通过百度搜索，魏则西和家人联系上北京武警二院的「生物免疫治疗中心」 \\
        2014.9 & 魏则西停止放化疗，花费重金进行免疫治疗。 \\
         & 付出大量医药费（超过20万）和时间后，癌细胞转移到肺，耽误了正规治疗。 \\
        2016.4 & 魏则西去世。 \\
        2016.5 & 联合调查组进驻百度公司，要求百度进行整改。 \\
        \end{tabular}
    \end{table}
\end{frame}


\begin{frame}
    \frametitle{}
    \begin{figure}[htpb]
        \centering
        \includegraphics[width=.8\linewidth]{改了吗}
    \end{figure}
\end{frame}


\begin{frame}
    \frametitle{}
    \begin{figure}[htpb]
        \centering
        \includegraphics[width=\linewidth]{百度问病}
    \end{figure}
\end{frame}


\section{信息安全}

\begin{frame}
    \frametitle{警惕高速下载陷阱}
    \begin{figure}[htpb]
        \centering
        \subfloat {\includegraphics[height=6.6cm]{search_baidu}}\qquad
        \subfloat {\includegraphics[height=6.6cm]{search_bing}} 
    \end{figure}
\end{frame}

\begin{frame}
    \frametitle{密码安全}
    \begin{minipage}{0.58\linewidth}
        \begin{figure}[htpb]
            \centering
            \includegraphics[width=\linewidth]{password}
        \end{figure}
    \end{minipage}\hfill
    \begin{minipage}{0.41\linewidth}
        账号密码安全建议：
        \begin{itemize}
          \item 保守口令密码，不告知他人
          \item 避免单纯使用个人信息设置密码
          \item 不使用弱密码，密码包含大写字母以及特殊字符
          \item 定期更换密码
          \item 不要所有系统使用同一个密码
        \end{itemize}
    \end{minipage}
\end{frame}

\begin{frame}
    \frametitle{明辨视听}
\begin{center}
    \cbf{国家网信办部署开展“清朗”系列专项行动}
\end{center}
    \qquad 2021年“清朗”系列专项行动从8个方面发力，包括整治网上历史虚无主义，治理算法滥用行为，打击\bf{网络水军}、流量造假、黑公关，整治未成年人网络环境，整治PUSH弹窗新闻信息突出问题，规范网站账号运营，整治网上文娱及热点排行乱象等。

    \qquad 2022年“清朗”系列专项行动将重点开展打击\bf{网络谣言}，打击网络直播、短视频领域乱象，算法综合治理，暑期未成年人网络环境整治、打击流量造假、黑公关、\bf{网络水军}等十方面任务。
    
    \qquad 2023年“清朗”系列专项行动，将聚焦新情况新问题以及制约治理成效的难点，开展一系列专项整治，其中9方面问题是重中之重，具体包括：\bf{整治“自媒体”乱象；打击“网络水军”操纵信息内容}；规范重点流量环节网络传播秩序；优化营商网络环境，保护企业合法权益；整治生活服务类平台信息内容乱象；整治短视频信息内容导向不良问题；整治暑期未成年人网络环境；整治网络戾气；整治春节网络环境。
\end{frame}


\begin{frame}
    \frametitle{网络实名制}
    \begin{center}
        \cbf{中华人民共和国网络安全法}

\small（2016年11月7日第十二届全国人民代表大会常务委员会第二十四次会议通过）

    \end{center}
    \cbf{第二十四条}\quad 网络运营者为用户办理网络接入、域名注册服务，办理固定电话、移动电话等入网手续，或者为用户提供信息发布、即时通讯等服务，在与用户签订协议或者确认提供服务时，\bf{应当要求用户提供真实身份信息。用户不提供真实身份信息的，网络运营者不得为其提供相关服务}。国家实施网络可信身份战略，支持研究开发安全、方便的电子身份认证技术，推动不同电子身份认证之间的互认。
\end{frame}

\begin{frame}
    \frametitle{网络实名制}
    \begin{center}
        \cbf{互联网用户账号信息管理规定}

\small（2022年6月9日国家互联网信息办公室2022年第11次室务会议审议通过）

    \end{center}
    \cbf{第十二条}\quad 互联网信息服务提供者应当在互联网用户账号信息页面展示合理范围内的互联网用户账号的互联网协议\bf{（IP）地址归属地}信息，便于公众为公共利益实施监督。

    \cbf{第十三条}\quad 互联网信息服务提供者应当在互联网用户公众账号信息页面，展示公众账号的运营主体、注册运营地址、内容生产类别、统一社会信用代码、有效联系方式、互联网协议\bf{（IP）地址归属地}等信息。
\end{frame}


\begin{frame}
    \frametitle{保护个人隐私}
    \begin{minipage}{0.48\linewidth}
        \begin{figure}[htpb]
            \centering
            \includegraphics[width=\linewidth]{private}
        \end{figure}
    \end{minipage}\hfill
    \begin{minipage}{0.48\linewidth}
        \bf{连选择权都没有，何谈“愿意”？}

        开启位置换来了更精准的外卖和打车服务；

        用阅读习惯换取了更精准的新闻推荐；

        用购买习惯获得了更精准的商品推送；

        用搜索习惯得到了更精准的内（guang）容（gao）推送……

        可是，这些隐私开启真的是我们自主愿意的吗？

        显然不是，事实上是我们“不得不”用自己的隐私获取便利，因为我们没得选择。
    \end{minipage}
\end{frame}

% 中国大多数互联网企业的崛起 ，都受益于排斥外来竞争者和不对称权责关系的基本盘，而互联网天然的马太效应，又强化了其垄断地位，进而巩固了前两个基本盘。

% 排斥外来竟争者，不必细说，这也是贸易战想要撼动的领域之一。

% 不对称权责关系。比如一边要求网络实名制 ，另一边却无对应的保护与惩处手段。一边要求用户让渡隐私，另一边却没有对应的制约手段。

% 以霸王条款获取用户信息，几乎存在于所有互联网公司。 至于订票平台 、 网络论坛泄露用户信息 ，也是家常便饭 ，最终顶多道歉了事。至于大数据杀熟 ，任你曝光 ， 我自岿然不动 。

% 再说互联网垄断。马太效应之下 ，互联网天然有滑向垄断的趋势 ，这倒不一定是坏事，规模经济和生态链优势反而能增加社会效益 。不过，当最终市场上只剩一两家寡头时，不仅消费者的选择权严重受限，而且我们在互联网上的几乎一切活动都受制于几家互联网企业。

% 有人说，国际上也是如此 ， 这是普遍现象。不过 ，缺乏限制的垄断与法律民意制约的垄断，完全不是一回事。 从 FACEBOOK 最近的遭遇就可以看得很清楚 。

% 遥想很多年前，互联网被视为打破利益阻碍的工具，没想到，随着互联网企业的壮大，它们已经巩固成为新的利益，和既有的利益一体，牢不可破。


\begin{frame}
    \frametitle{保护个人隐私}
    \begin{center}
        \cbf{工业和信息化部关于侵害用户权益行为的APP通报（2023年第1批，总第27批）}
    \end{center}
\vspace{-1em}
\begin{table}[]
    \begin{tabular}{ccl}
    \rowcolor[HTML]{007AAA} 
    \multicolumn{1}{l}{\cellcolor[HTML]{007AAA}{\color[HTML]{FFFFFF} \textbf{序号}}} & \multicolumn{1}{l}{\cellcolor[HTML]{007AAA}{\color[HTML]{FFFFFF} \textbf{应用名称}}} & {\color[HTML]{FFFFFF} \textbf{所涉问题}}            \\
    \rowcolor[HTML]{FFFFFF} 
    {\color[HTML]{070707} 1}                                                       & {\color[HTML]{070707} 墨迹天气极速版}                                                   & {\color[HTML]{070707} 应用分发平台上的APP信息明示不到位}       \\
    \rowcolor[HTML]{EDF5FA} 
    {\color[HTML]{070707} 2}                                                       & {\color[HTML]{070707} 掌上公交}                                                      & {\color[HTML]{070707} 违规互联网弹窗信息推送服务}            \\
    \rowcolor[HTML]{FFFFFF} 
    {\color[HTML]{070707} 3}                                                       & {\color[HTML]{070707} 华为阅读}                                                      & {\color[HTML]{070707} 欺骗误导强迫用户}                 \\
    \rowcolor[HTML]{EDF5FA} 
    {\color[HTML]{070707} 4}                                                       & {\color[HTML]{070707} 米读极速版}                                                     & {\color[HTML]{070707} 欺骗误导强迫用户}                 \\
    \rowcolor[HTML]{FFFFFF} 
    {\color[HTML]{070707} 5}                                                       & {\color[HTML]{070707} 消灭星星全新版}                                                   & {\color[HTML]{070707} 违规互联网弹窗信息推送服务}            \\
    \rowcolor[HTML]{EDF5FA} 
    {\color[HTML]{070707} 6}                                                       & {\color[HTML]{070707} 小初高同步课堂}                                                   & {\color[HTML]{070707} 超范围收集个人信息}                \\
    \rowcolor[HTML]{FFFFFF} 
    {\color[HTML]{070707} 7}                                                       & {\color[HTML]{070707} iGola骑鹅旅行}                                                 & {\color[HTML]{070707} 违规收集、使用个人信息；APP强制、频繁、过度索取权限}        \\
    \rowcolor[HTML]{EDF5FA} 
    {\color[HTML]{070707} 8}                                                       & {\color[HTML]{070707} 导航卫星地图}                                                    & {\color[HTML]{070707} 违规收集个人信息；APP强制、频繁、过度索取权限}\\
    \rowcolor[HTML]{FFFFFF} 
    {\color[HTML]{070707} 9}                                                       & {\color[HTML]{070707} 醉美青海}                                                    & {\color[HTML]{070707} 违规收集、使用个人信息；APP强制、频繁、过度索取权限}\\
    \rowcolor[HTML]{EDF5FA} 
    {\color[HTML]{070707} 10}                                                       & {\color[HTML]{070707} 本地出行}                                                    & {\color[HTML]{070707} 违规收集个人信息；APP强制、频繁、过度索取权限}
    \end{tabular}
    \end{table}
\end{frame}

\section{一些基础操作}


\begin{frame}
    \frametitle{切换输入法}
    \begin{itemize}
      \item 切换键盘布局 Ctrl + Shift
      \item 打开或关闭中文输入法编辑器 Ctrl + Space
      \item 切换输入语言 Alt + Shift
    \end{itemize}
\end{frame}


\begin{frame}
    \frametitle{截屏}
    \begin{minipage}{0.33\linewidth}
        \begin{figure}[htpb]
            \centering
            \includegraphics[width=\linewidth]{printscreen}
        \end{figure}
    \end{minipage}\hfill
    \begin{minipage}{0.63\linewidth}
        \begin{itemize}
          \item 全屏截图：PrintScreen
          \item QQ截图：Ctrl + Shift + A
          \item 微信截图：Alt + A
          \item Windows10截屏功能：Win + Shift + S
        \end{itemize}
    \end{minipage}
\end{frame}


\begin{frame}
    \frametitle{文件管理}
    \begin{enumerate}
      \item 打开资源管理器
      \item 资源管理器的界面布局
      \item 文件名和文件类型
      \item 文档属性（只读/隐藏）
      \item 显示文件扩展名、显示隐藏文件
      \item 查找文件
      \item 文件和文件夹的创建、复制、粘贴、移动、删除
      \item 文件路径
      \item 建立快捷方式
    \end{enumerate}
\end{frame}

\begin{frame}
    \frametitle{任务管理器}
    \begin{itemize}
      \item 查看当前运行程序列表
      \item 查看硬件性能消耗
      \item 结束进程
    \end{itemize}
\end{frame}




\end{document}
